\documentclass[12pt,a4paper]{report}

\usepackage[spanish,es-tabla]{babel}
\usepackage[utf8]{inputenc}
\usepackage[T1]{fontenc}
\usepackage{helvet}
\renewcommand{\familydefault}{\sfdefault}

\usepackage[
  top=3cm, bottom=3cm, left=3cm, right=3cm,
  headheight=15pt, headsep=10pt
]{geometry}
\usepackage{setspace}
\setstretch{1.15}
\setlength{\parindent}{1.25cm}
\setlength{\parskip}{0.5em}

\usepackage[table]{xcolor}
% Paleta alineada con frontend/css/dr-agramonte.css
\definecolor{teal}{RGB}{15,118,110}
\definecolor{tealdark}{RGB}{19,78,72}
\definecolor{teallight}{RGB}{20,184,166}
\definecolor{tealpale}{RGB}{240,253,250}
\definecolor{silver}{RGB}{212,212,216}
\definecolor{offwhite}{RGB}{249,250,251}
\definecolor{graytext}{RGB}{107,114,128}
\definecolor{graydark}{RGB}{31,41,55}
\definecolor{carddark}{RGB}{30,41,59}
\definecolor{bgdark}{RGB}{15,23,42}
\definecolor{rowgray}{RGB}{243,244,246}
\definecolor{bordergray}{RGB}{229,231,235}
\definecolor{codebg}{RGB}{30,41,59}
\definecolor{codefg}{RGB}{226,232,240}

% Alias compatibles con el estilo anterior
\colorlet{darkblue}{tealdark}
\colorlet{medblue}{teal}
\colorlet{lightblue}{tealpale}

\tikzset{
  block/.style={
    draw=teal, rounded corners=10pt, line width=0.9pt,
    fill=tealpale, text=graydark, font=\small\sffamily,
    minimum height=1.1cm, align=center, inner sep=6pt
  },
  blockdark/.style={
    block, fill=carddark, text=white, draw=teallight
  },
  blockaccent/.style={
    block, fill=teal!85!black, text=white, draw=tealdark
  },
  arrow/.style={->, >=Stealth, line width=0.9pt, teal},
  entity/.style={
    draw=teal, fill=tealpale, rounded corners=6pt,
    font=\small\sffamily, align=center, inner sep=5pt
  },
  rel/.style={->, >=Stealth, teal, line width=0.7pt}
}


\usepackage{amssymb}
\usepackage{pifont}
\DeclareUnicodeCharacter{2713}{\checkmark}

\usepackage{fancyhdr}
\pagestyle{fancy}
\fancyhf{}
\fancyhead[L]{\small\color{tealdark}Mar Agramonte --- Proyecto DAM}
\fancyhead[R]{\small\thepage}
\renewcommand{\headrulewidth}{0.4pt}

\usepackage{titlesec}
\titleformat{\chapter}[display]
  {\normalfont\Large\bfseries\color{tealdark}}
  {Capítulo \thechapter:}{16pt}{\Large}
\titleformat{\section}{\normalfont\large\bfseries\color{tealdark}}{\thesection}{1em}{}
\titleformat{\subsection}{\normalfont\normalsize\bfseries\color{teal}}{\thesubsection}{1em}{}

\usepackage{tocloft}
\renewcommand{\cftchapfont}{\bfseries\color{tealdark}}
\renewcommand{\cftchappagefont}{\bfseries\color{tealdark}}

\usepackage{booktabs, array, longtable, float, caption}
\captionsetup{font=small,labelfont={bf,color=tealdark}}

\usepackage{listings}
\lstdefinestyle{mycode}{
  backgroundcolor=\color{codebg},
  basicstyle=\ttfamily\small\color{codefg},
  breaklines=true, numbers=left, numbersep=8pt,
  frame=single, framerule=0pt, xleftmargin=14pt,
}
\lstset{style=mycode}

\usepackage{graphicx}
\usepackage{tikz}
\usetikzlibrary{positioning, arrows.meta, shapes.geometric}

\usepackage{hyperref}
\hypersetup{
  colorlinks=true, linkcolor=tealdark, urlcolor=teal, citecolor=tealdark,
  pdftitle={Dr. Agramonte - Gestión de Citas Médicas},
  pdfauthor={Mar Agramonte},
}

\newcommand{\hrow}{\rowcolor{tealpale}}
\newcommand{\grayrow}{\rowcolor{rowgray}}

\begin{document}

% --- Portada ---
\begin{titlepage}
  \pagestyle{empty}
  \centering
  \vspace*{1.5cm}
  {\Large\bfseries\color{tealdark} TRABAJO DE FIN DE CICLO\\[0.8em]}
  \vspace{0.8cm}
  {\Huge\bfseries\color{teal} DR. AGRAMONTE\\[0.4em]}
  {\Large\bfseries\color{teallight} Plataforma web de gestión de citas médicas\\[1.5em]}
  \begin{tikzpicture}
    \fill[tealpale, rounded corners=12pt] (0,0) rectangle (12,0.35);
    \fill[teal] (0,0) rectangle (4,0.35);
  \end{tikzpicture}
  \vfill
  \begin{tabular}{rl}
    \textbf{Autora:} & Mar Agramonte \\
    \textbf{Centro:} & CESUR \\
    \textbf{Ciclo:} & Desarrollo de Aplicaciones Multiplataforma \\
    \textbf{Tutores:} & Pascual Barrer -- Jaume Rossinrol \\
    \textbf{Curso:} & 2024--2025 \\
    \textbf{Fecha:} & Mayo 2026 \\
  \end{tabular}
\end{titlepage}

\clearpage
\tableofcontents
\newpage

% ============================================================================
\chapter{Resumen ejecutivo}
% ============================================================================

\textbf{Dr. Agramonte} es una aplicación web para que pacientes reserven, consulten y cancelen citas médicas sin depender del teléfono del profesional. El sistema integra frontend estático (HTML, CSS y JavaScript), API REST con \textbf{Spring Boot 3.5.3} y \textbf{PostgreSQL}, autenticación \textbf{JWT}, control de concurrencia con \textbf{bloqueo pesimista}, notificaciones \textbf{Twilio} (SMS/WhatsApp) y despliegue con \textbf{Docker Compose}.

La interfaz replica la identidad visual del consultorio: paleta \textbf{teal} (\#0F766E, \#14B8A6), tipografías DM Sans / DM Serif Display, modo claro/oscuro y diseño responsive. La reserva se organiza en pasos (centro, calendario, hora, datos y confirmación) y puede iniciarse desde \texttt{servicios.html} con el tipo de consulta preseleccionado.

\textbf{Alcance real del repositorio} (mayo 2026): no incluye MongoDB, Redis ni microservicio Node.js; la agenda médico y las pruebas de reserva se consultan en \texttt{panel-pruebas.html}. Las mejoras futuras se documentan como evolución, no como funcionalidad entregada.

% ============================================================================
\chapter{Justificación y contexto}
% ============================================================================

\section{Origen del problema}

La motivación proviene de la observación directa del sector sanitario: gran parte de las llamadas al móvil del médico no son urgencias clínicas, sino gestión de agenda. Una web accesible reduce interrupciones y devuelve tiempo al profesional.

\section{Competidores y hueco de mercado}

Plataformas como Doctoralia o Top Doctors cubren grandes centros, pero implican comisiones y dependencia de terceros. Consultas pequeñas o médicos autónomos necesitan una herramienta simple, bajo su control y adaptable a medicina interna.

\section{Marco legal}

Los datos de salud exigen cumplir \textbf{RGPD}, \textbf{LOPDGDD} y \textbf{Ley 41/2002}. La aplicación incluye política de privacidad, consentimiento en reserva y minimización de datos (solo los necesarios para la cita).

% ============================================================================
\chapter{Objetivos}
% ============================================================================

\section{Objetivo general}

Desarrollar una aplicación web de gestión de citas médicas que permita al paciente reservar y cancelar de forma autónoma, con backend seguro y despliegue reproducible.

\section{Objetivos específicos (alineados con la implementación)}

\begin{enumerate}
  \item Modelar y persistir usuarios, médicos, horarios y citas en \textbf{PostgreSQL} con migraciones \textbf{Flyway} y perfiles \texttt{postgres}/\texttt{mysql}.
  \item Implementar API REST con \textbf{Spring Boot}, autenticación \textbf{JWT} y roles \texttt{PACIENTE}, \texttt{MEDICO} y \texttt{ADMIN} en \texttt{SecurityConfig}.
  \item Garantizar integridad de reservas con \texttt{SELECT FOR UPDATE} sobre horarios disponibles.
  \item Construir frontend responsive, PWA, modo oscuro y flujo de reserva accesible (\texttt{reserva.js}, \texttt{api-client.js}).
  \item Enviar notificaciones al médico y al paciente vía \textbf{Twilio}, con recordatorio programado 24\,h antes (\texttt{CitaReminderScheduler}).
  \item Desplegar con \textbf{Docker Compose} (nginx + backend + base de datos).
  \item Documentar el proyecto y preparar defensa con demo reproducible.
\end{enumerate}

% ============================================================================
\chapter{Análisis y contextualización del sector}
% ============================================================================

\section{Sector profesional}

El ciclo DAM aporta competencias en desarrollo multiplataforma, bases de datos, interfaces y despliegue. Este proyecto las aplica al ámbito sanitario privado: agenda, relación paciente--médico y comunicación.

\section{Riesgos laborales (resumen)}

\begin{itemize}
  \item \textbf{Ergonómicos:} posturas prolongadas ante pantalla; pausas y ergonomía en puesto de desarrollo y en consulta.
  \item \textbf{Psicosociales:} estrés por interrupciones; la digitalización de citas reduce carga administrativa.
  \item \textbf{Seguridad de la información:} tratamiento de datos de salud; cifrado en tránsito (HTTPS), contraseñas con BCrypt, secretos en variables de entorno.
  \item \textbf{EPI / buenas prácticas:} control de accesos por rol, copias de seguridad de BD, no commitear credenciales.
\end{itemize}

% ============================================================================
\chapter{Diseño y arquitectura del sistema}
% ============================================================================

\section{Arquitectura general}

La solución sigue una arquitectura en capas desplegada con contenedores. El navegador consume el frontend servido por \textbf{nginx}; las peticiones \texttt{/api} se proxyan al backend Java. Twilio es un servicio externo invocado desde \texttt{TwilioMessageService}.

\begin{figure}[H]
\centering
\begin{tikzpicture}[node distance=1.2cm and 1.4cm, font=\small\sffamily]
  \node[block, minimum width=2.6cm] (client) {Navegador\\Paciente / Médico};
  \node[block, minimum width=2.4cm, right=of client] (nginx) {Nginx\\Frontend estático\\+ proxy \texttt{/api}};
  \node[blockaccent, minimum width=2.8cm, right=of nginx] (api) {Spring Boot 3.5\\API REST + JWT\\Java 21};
  \node[block, minimum width=2.4cm, below=of api] (db) {PostgreSQL 15\\Flyway + JPA};
  \node[block, minimum width=2.2cm, above right=0.9cm and 1.1cm of api] (twilio) {Twilio\\SMS / WhatsApp};

  \draw[arrow] (client) -- node[above, font=\scriptsize, text=graytext]{HTTPS :80} (nginx);
  \draw[arrow] (nginx) -- node[above, font=\scriptsize, text=graytext]{/api} (api);
  \draw[arrow] (api) -- node[right, font=\scriptsize, text=graytext]{JDBC} (db);
  \draw[arrow] (api) -- node[above, font=\scriptsize, text=graytext]{REST API} (twilio);

  \node[below=0.15cm of client, font=\scriptsize, text=graytext] {PWA + \texttt{service-worker.js}};
  \node[below=0.15cm of nginx, font=\scriptsize, text=graytext] {HTML/CSS/JS};
  \node[below=0.15cm of db, font=\scriptsize, text=graytext] {Perfil \texttt{postgres} o \texttt{mysql}};
\end{tikzpicture}
\caption{Arquitectura desplegada con Docker Compose (implementación actual)}
\label{fig:arquitectura}
\end{figure}


\section{Stack tecnológico}

\begin{table}[H]
\centering
\caption{Componentes implementados}
\begin{tabular}{>{\bfseries\color{white}\columncolor{teal}}l >{\columncolor{white}}p{4.2cm} >{\columncolor{white}}p{5.5cm}}
\toprule
\color{white}Capa & \color{white}Tecnología & \color{white}Función \\
\midrule
Frontend & HTML5, CSS3, JS ES6+ & UI, PWA, \texttt{reserva.js} \\
\rowcolor{rowgray}
Proxy / estático & nginx (Docker) & Archivos + proxy \texttt{/api} \\
Backend & Java 21, Spring Boot 3.5.3 & API, seguridad, negocio \\
\rowcolor{rowgray}
Persistencia & PostgreSQL 15, Flyway, JPA & Datos transaccionales \\
Notificaciones & Twilio API (Java SDK) & SMS / WhatsApp \\
\rowcolor{rowgray}
Despliegue & Docker Compose & Entorno reproducible \\
\bottomrule
\end{tabular}
\end{table}

\section{Modelo de datos}

\begin{figure}[H]
\centering
\begin{tikzpicture}[font=\small\sffamily, node distance=0.8cm and 1.6cm]
  \node[entity, minimum width=2.4cm] (u) {\textbf{usuarios}\\\scriptsize id, email, password\\nombre, teléfono, rol};
  \node[entity, minimum width=2.2cm, right=of u] (m) {\textbf{medicos}\\\scriptsize id, nombre\\especialidad, email};
  \node[entity, minimum width=2.5cm, below=1.1cm of m] (h) {\textbf{horarios}\\\scriptsize medico\_id, inicio, fin\\disponible};
  \node[entity, minimum width=2.6cm, below=1.1cm of u] (c) {\textbf{citas}\\\scriptsize usuario\_id, medico\_id, centro\_id\\fecha\_hora, motivo, estado\\recordatorio\_enviado\\cobertura, pago (V8)};
  \node[entity, minimum width=2.5cm, right=2.2cm of c] (n) {\textbf{notificaciones}\\\scriptsize usuario\_id, cita\_id\\canal, estado\_entrega};
  \node[entity, minimum width=2.4cm, below=1.0cm of c] (ce) {\textbf{centros}\\\scriptsize codigo, nombre\\direccion, ciudad};

  \draw[rel] (u) -- node[above, sloped, font=\scriptsize]{1:N} (c);
  \draw[rel] (m) -- node[above, sloped, font=\scriptsize]{1:N} (c);
  \draw[rel] (m) -- node[right, font=\scriptsize]{1:N} (h);
  \draw[rel, dashed] (c) -- node[above, font=\scriptsize]{opc.} (n);
  \draw[rel, dashed] (u) -- (n);
  \draw[rel, dashed] (ce) -- node[right, font=\scriptsize]{opc. 1:N} (c);
  \draw[rel, dashed] (ce) to[bend left=20] node[left, font=\scriptsize]{opc. 1:N} (h);
\end{tikzpicture}
\caption{Modelo entidad-relación principal (PostgreSQL, migraciones Flyway V1--V8)}
\label{fig:er}
\end{figure}


La migración \texttt{V2\_\_cita\_recordatorio} añade el campo \texttt{recordatorio\_enviado} para evitar duplicar avisos Twilio.

\section{Frontend y paleta visual}

Los colores del documento y de los diagramas coinciden con \texttt{dr-agramonte.css}: teal principal, fondos \texttt{off-white}, tarjetas blancas y variante oscura (\texttt{data-theme="dark"}, fondo \#0F172A).

\begin{figure}[H]
\centering
\begin{tikzpicture}[node distance=0.7cm, font=\small\sffamily]
  \node[blockaccent, minimum width=3.2cm] (pages) {Páginas HTML\\index, servicios, reservar...};
  \node[block, minimum width=3cm, below=of pages] (css) {\texttt{dr-agramonte.css}\\Tema claro/oscuro\\\texttt{data-theme}};
  \node[block, minimum width=2.8cm, below left=0.6cm and 0.4cm of css] (js) {\texttt{reserva.js}\\\texttt{api-client.js}};
  \node[block, minimum width=2.6cm, below right=0.6cm and 0.4cm of css] (pwa) {Service Worker\\modo offline};

  \draw[arrow] (pages) -- (css);
  \draw[arrow] (css) -- (js);
  \draw[arrow] (css) -- (pwa);
  \node[right=0.3cm of js, font=\scriptsize, text=graytext, align=left] {Centro $\rightarrow$ calendario\\$\rightarrow$ slots API $\rightarrow$ confirmar};
\end{tikzpicture}
\caption{Capas del frontend y flujo de reserva en cuatro pasos}
\label{fig:frontend-capas}
\end{figure}


% ============================================================================
\chapter{Desarrollo técnico}
% ============================================================================

\section{Backend}

\subsection{Autenticación JWT}

\texttt{AuthController} expone \texttt{POST /api/auth/login} y \texttt{/api/auth/registro}. El token se envía en cabecera \texttt{Authorization: Bearer}. \texttt{JwtAuthFilter} valida cada petición protegida.

\subsection{Reserva y concurrencia}

\begin{lstlisting}[language=Java]
Horario horario = horarioRepository
  .findByMedicoIdAndInicioAndDisponibleTrueForUpdate(...)
  .orElseThrow(...);
\end{lstlisting}

Tras validar, se crea la \texttt{Cita}, se marca el horario no disponible y se invoca \texttt{CitaNotificationService}.

\begin{figure}[H]
\centering
\begin{tikzpicture}[node distance=0.55cm, font=\small\sffamily,
  every node/.style={align=center}]
  \node[draw=teal, fill=tealpale, rounded corners=6pt, inner sep=5pt] (fe) {Frontend\\\texttt{reserva.js}};
  \node[draw=teal, fill=tealpale, rounded corners=6pt, inner sep=5pt, below=of fe] (api) {POST \texttt{/api/citas}\\JWT paciente};
  \node[draw=teal, fill=tealpale, rounded corners=6pt, inner sep=5pt, below=of api] (svc) {\texttt{CitaService.crearCita}};
  \node[draw=tealdark, fill=teal!20, rounded corners=6pt, inner sep=5pt, below=of svc] (lock) {SELECT ... FOR UPDATE\\horario disponible};
  \node[draw=teal, fill=tealpale, rounded corners=6pt, inner sep=5pt, below=of lock] (save) {INSERT cita + UPDATE horario};
  \node[draw=teallight, fill=teal!30, text=white, rounded corners=6pt, inner sep=5pt, below=of save] (tw) {Twilio: médico + paciente};

  \draw[arrow] (fe) -- (api);
  \draw[arrow] (api) -- (svc);
  \draw[arrow] (svc) -- (lock);
  \draw[arrow] (lock) -- (save);
  \draw[arrow] (save) -- (tw);
\end{tikzpicture}
\caption{Flujo de reserva de cita con bloqueo pesimista y notificaciones}
\label{fig:secuencia-reserva}
\end{figure}


\subsection{API principal}

\begin{table}[H]
\centering
\small
\begin{tabular}{llp{6.2cm}}
\toprule
\textbf{Método} & \textbf{Ruta} & \textbf{Descripción} \\
\midrule
POST & /api/auth/login, /registro & Sesión JWT \\
GET & /api/medicos & Listado de especialistas \\
GET & /api/disponibilidad & Slots por médico y fecha \\
POST & /api/citas & Crear cita (paciente autenticado) \\
GET & /api/citas/mias & Citas del paciente \\
DELETE & /api/citas/\{id\} & Cancelación \\
POST & /api/contacto & Formulario de contacto \\
\bottomrule
\end{tabular}
\caption{Endpoints principales}
\end{table}

\section{Notificaciones Twilio}

\texttt{CitaNotificationService} envía:
\begin{itemize}
  \item Aviso al médico (\texttt{TWILIO\_NOTIFY\_TO}) en alta y cancelación.
  \item Confirmación y cancelación al paciente si tiene teléfono.
  \item Recordatorio $\sim$24\,h antes vía \texttt{@Scheduled} (ventana 23--25\,h).
\end{itemize}

Variables: \texttt{TWILIO\_ENABLED}, \texttt{TWILIO\_ACCOUNT\_SID}, \texttt{TWILIO\_AUTH\_TOKEN}, \texttt{TWILIO\_FROM}, \texttt{TWILIO\_CHANNEL} (whatsapp o sms).

\section{Frontend}

\begin{itemize}
  \item \textbf{Reserva:} \texttt{reservar.html} + \texttt{reserva.js} (centros, calendario, slots desde API, fallback local).
  \item \textbf{Servicios:} enlaces \texttt{reservar.html?servicio=...} preconfiguran tipo y motivo.
  \item \textbf{Tema:} botón luna/sol, clase \texttt{btn--reserva} unificada en todo el sitio.
  \item \textbf{PWA:} \texttt{manifest.json} y \texttt{service-worker.js}; página \texttt{offline.html}.
\end{itemize}

\section{Seguridad y RGPD}

\begin{itemize}
  \item Contraseñas con \texttt{BCryptPasswordEncoder(12)}.
  \item CORS configurado para orígenes de desarrollo y producción.
  \item Política de privacidad en \texttt{privacidad.html}.
  \item Checkbox de aceptación en formulario de reserva.
  \item Secretos (\texttt{JWT\_SECRET}, Twilio) en \texttt{.env}, no en el repositorio.
\end{itemize}

\section{Despliegue}

\begin{lstlisting}[language=bash]
docker compose up -d --build
\end{lstlisting}

Servicios: \texttt{postgres}, \texttt{backend} (:8080), \texttt{frontend} (:80). Perfil opcional MySQL con \texttt{--profile mysql}.

% ============================================================================
\chapter{Diagramas y modelos}
% ============================================================================

Los diagramas de este capítulo usan la paleta teal del CSS del proyecto. Para actualizarlos, editar los ficheros \texttt{docs/tikz-diagrama-*.tex}.

% ============================================================================
\chapter{Planificación temporal}
% ============================================================================

\begin{center}
\begin{tikzpicture}[x=0.45cm, y=0.55cm, font=\scriptsize\sffamily]
  \draw[very thin, bordergray] (0,0) grid (24,8);
  \foreach \x/\txt in {2/M1, 8/M2, 14/M3, 20/M4}
    \node at (\x+1,8.6) [font=\bfseries, text=tealdark] {\txt};
  \node[left, text=tealdark] at (0,7) {Análisis};
  \node[left, text=tealdark] at (0,6) {Diseño BD/API};
  \node[left, text=tealdark] at (0,5) {Backend};
  \node[left, text=tealdark] at (0,4) {Frontend};
  \node[left, text=tealdark] at (0,3) {Twilio + pruebas};
  \node[left, text=tealdark] at (0,2) {Docker};
  \node[left, text=tealdark] at (0,1) {Memoria};
  \fill[teal] (0,6.7) rectangle (4,7.3);
  \fill[tealdark] (3,5.7) rectangle (7,6.3);
  \fill[teal] (6,4.7) rectangle (13,5.3);
  \fill[teallight] (12,3.7) rectangle (18,4.3);
  \fill[teal] (11,2.7) rectangle (16,3.3);
  \fill[tealdark] (17,1.7) rectangle (21,2.3);
  \fill[tealpale, draw=teal] (18,0.7) rectangle (24,1.3);
\end{tikzpicture}
\end{center}

% ============================================================================
\chapter{Pruebas y validación}
% ============================================================================

\section{Pruebas realizadas}

\begin{itemize}
  \item Compilación Maven (\texttt{mvn -DskipTests compile}).
  \item Arranque Docker Compose y health check API.
  \item Flujo manual: registro/login $\rightarrow$ reservar $\rightarrow$ mis citas $\rightarrow$ cancelar.
  \item Prueba de parámetro \texttt{?servicio=} desde servicios.
  \item Validación de formularios (email, teléfono español).
\end{itemize}

\section{Pruebas pendientes (mejora futura)}

\begin{itemize}
  \item Suite JUnit + Testcontainers (dependencias ya en \texttt{pom.xml}).
  \item Tests E2E (Playwright/Cypress).
  \item Auditoría Lighthouse formal documentada.
\end{itemize}

% ============================================================================
\chapter{Conclusiones}
% ============================================================================

Se ha entregado una aplicación funcional y desplegable que responde al problema de gestión de citas, con arquitectura clara y código mantenible. Los objetivos técnicos del repositorio se cumplen en PostgreSQL, Spring Boot, frontend PWA y Twilio integrado en Java.

\textbf{Limitaciones honestas:} panel médico sin API de historial, sin refresh token, sin tests automatizados completos. Son líneas de evolución natural del TFG.

\textbf{Reflexión:} la coherencia entre memoria, código y demo oral es tan importante como la funcionalidad; este documento describe la versión real del sistema a mayo de 2026.

% ============================================================================
\chapter{Agradecimientos}
% ============================================================================

A las personas del entorno sanitario que inspiraron el problema, a los tutores del ciclo, a quienes probaron la reserva en distintos dispositivos, y a mi familia por el apoyo durante el desarrollo.

% ============================================================================
\chapter{Checklist de cumplimiento (rúbricas)}
% ============================================================================

\enlargethispage{2\baselineskip}
\begin{center}
  {\large\bfseries\color{tealdark} Checklist de cumplimiento --- Módulo de Proyecto}\\[0.3em]
  {\small Proceso (25\,\%) + contenido escrito (40\,\%) + defensa (35\,\%)}\\[0.2em]
  {\footnotesize Memoria alineada con el repositorio (mayo 2026). Marcar al entregar.}
\end{center}

\vspace{0.3em}
\scriptsize
\setlength{\tabcolsep}{3pt}
\renewcommand{\arraystretch}{1.12}

\begin{longtable}{@{}p{0.34\textwidth}ccp{0.36\textwidth}@{}}
\toprule
\textbf{Criterio} & \textbf{Mem.} & \textbf{Demo} & \textbf{Notas} \\
\midrule
\endfirsthead
\toprule
\textbf{Criterio} & \textbf{Mem.} & \textbf{Demo} & \textbf{Notas} \\
\midrule
\endhead

\multicolumn{4}{l}{\textbf{\color{tealdark}A. Proceso (25\,\%)}} \\
\midrule
Seguimiento con tutor documentado & $\square$ & --- & Anexar evidencias \\
Gantt / hitos coherentes con el calendario real & $\checkmark$ & --- & Cap. planificación \\
Coherencia memoria $\leftrightarrow$ código & $\checkmark$ & $\checkmark$ & Sin Mongo/Redis/Node en memoria \\
\midrule

\multicolumn{4}{l}{\textbf{\color{tealdark}B. Contenido escrito (40\,\%)}} \\
\midrule
Estructura guía (justificación, objetivos, sector, desarrollo, conclusiones) & $\checkmark$ & --- & \\
PDF, márgenes, fuente, extensión & $\square$ & --- & Revisar 35--65 páginas \\
Diagramas TikZ (paleta teal CSS) & $\checkmark$ & --- & Cap. arquitectura \\
Referencias APA & $\square$ & --- & Ampliar bibliografía si tutor exige \\
Anexos (capturas, manual usuario) & $\square$ & --- & Recomendado \\
\midrule

\multicolumn{4}{l}{\textbf{\color{tealdark}C. Defensa (35\,\%)}} \\
\midrule
Demo: login $\rightarrow$ reservar $\rightarrow$ cancelar & $\square$ & $\square$ & \texttt{docker compose up} \\
Demo servicio preseleccionado + modo oscuro & $\square$ & $\square$ & \\
Plan B (capturas / vídeo) & $\square$ & $\square$ & \\
Twilio en sandbox (opcional) & $\square$ & $\square$ & \\
\midrule

\multicolumn{4}{l}{\textbf{\color{tealdark}D. Implementación verificada en repo}} \\
\midrule
PostgreSQL + Flyway + Spring Boot + JWT & $\checkmark$ & $\checkmark$ & \\
\texttt{SELECT FOR UPDATE} + citas & $\checkmark$ & $\checkmark$ & \\
Frontend PWA + reserva + servicios URL & $\checkmark$ & $\checkmark$ & \\
Twilio Java (médico, paciente, recordatorio) & $\checkmark$ & $\square$ & Requiere \texttt{.env} \\
Panel médico API historial & $\square$ & $\square$ & Solo demo estática \\
Tests automatizados JUnit/E2E & $\square$ & $\square$ & Mejora futura \\
\bottomrule
\end{longtable}

\normalsize


% ============================================================================
\chapter{Bibliografía}
% ============================================================================

\begin{thebibliography}{99}

\bibitem{spring}
Spring. (2024). \textit{Spring Boot 3.5 Reference}. \url{https://spring.io/projects/spring-boot}

\bibitem{rgpd}
Reglamento (UE) 2016/679 (RGPD).

\bibitem{ley41}
Ley 41/2002, autonomía del paciente e historia clínica (BOE).

\bibitem{wcag}
W3C (2021). \textit{WCAG 2.1}. \url{https://www.w3.org/WAI/WCAG21/quickref/}

\bibitem{docker}
Docker. \textit{Compose Documentation}. \url{https://docs.docker.com/compose/}

\bibitem{twilio}
Twilio. \textit{Messaging API}. \url{https://www.twilio.com/docs/messaging}

\bibitem{flyway}
Flyway. \textit{Documentation}. \url{https://documentation.red-gate.com/fd}

\bibitem{doctoralia}
Doctoralia. \textit{Plataforma para profesionales}. \url{https://www.doctoralia.es}

\end{thebibliography}

\end{document}
